\section{Related Work}
\label{sec:related}

Prompt extraction and manipulation attacks against LLM applications have received considerable attention. PLeak~\cite{hui2024pleak} casts prompt extraction as an optimization problem and reconstructs system prompts almost completely from crafted queries; multi-turn leakage attacks~\cite{multiturn2024} exploit the sycophancy of instruction-tuned models to recover hidden instructions incrementally, with success rates that reach 86.2\%; and systematic reviews of prompt injection~\cite{promptinjection2026} show that keeping the prompt confidential is not sufficient to guarantee secure behavior.

Software-based defenses respond to these threats at specific points of the interaction pipeline. ProxyPrompt~\cite{proxyprompt2025} interposes a proxy that shields the real prompt, SysVec~\cite{sysvec2025} replaces the textual prompt with a non-invertible vector representation, PromptGuard~\cite{promptguard2026} applies structured filtering against injection, Encrypted Prompt~\cite{encryptedprompt2025} gates prompt access through permission-bearing tokens, PromptKeeper~\cite{promptkeeper2024} frames leakage detection as hypothesis testing, and Attention Tracker~\cite{attentiontracker2025} inspects model internals for attention patterns typical of adversarial inputs. Each operates inside a single trust boundary and targets one attack vector at a time. A second line of work composes several defenses: TRYLOCK~\cite{trylock2026} stacks jailbreak defenses and Countermind~\cite{countermind2025} adds auditing and monitoring to the inference pipeline, while analyses of serverless machine learning security~\cite{sas2026} describe the enlarged attack surface of cloud-native deployments; these contributions concentrate on inference-time robustness and leave cross-domain deployment unaddressed. At the infrastructure level, TEE-based systems such as Petridish~\cite{petridish2024} and Omega~\cite{omega2025} run inference inside secure enclaves, with moderate overheads under controlled conditions~\cite{teeperf2025}, but depend on specific hardware and offer no forensic attribution or cross-system monitoring. Watermarking addresses the complementary problem of output provenance: token-level schemes~\cite{kirchenbauer2023watermark}, scalable frameworks such as Waterfall~\cite{waterfall2024} and SynthID~\cite{synthid2024}, and zero-width character (ZWC) steganography~\cite{zheng2020zwc}, whose covert channel can also be abused~\cite{stegocollusion2024}; these techniques target attribution, not the confidentiality of system prompts.

Taken together, the literature treats prompt confidentiality, input validation, output filtering, and attribution as separate problems, usually inside one trust boundary, and proposals that combine several mechanisms rarely report how much each component contributes or whether their effectiveness holds when the underlying model changes. This article composes cryptographic protection, validation, filtering, watermarking, and monitoring across layers and administrative boundaries, and evaluates the result across model families and layer configurations.
